\section{Critical Groups and the 2-Adic Transfer Law}
\label{sec:smith}

\subsection{The Smith Normal Form Invariant}

Let $G$ be a graph with adjacency matrix $A$. The \emph{critical group} (or
\emph{sandpile group}) of $G$ is the torsion part of $\mathbb{Z}^n / \mathrm{Im}(A)$,
determined up to isomorphism by the Smith normal form (SNF) of $A$.
Write $\mathrm{SNF}(A) = \mathrm{diag}(s_1, \ldots, s_n)$ with $s_i \mid s_{i+1}$.

For a strongly regular graph $\mathrm{SRG}(v,k,\lambda,\mu)$ the adjacency spectrum
is $\{k^1, r^f, s^g\}$ and $\det(A) = k \cdot r^f \cdot s^g$ is a complete spectral invariant.

\subsection{The 2-Adic Transfer Theorem}

\begin{theorem}[2-Adic Transfer Law]
\label{thm:transfer}
Let $G_1, G_2$ be cospectral graphs (same adjacency eigenvalue multiset).
For every prime $p$:
$$\sum_{i=1}^{n} v_p(s_i(G_1)) = \sum_{i=1}^{n} v_p(s_i(G_2)).$$
That is, the total $p$-adic valuation of the Smith group is a spectral invariant.
\end{theorem}

\begin{proof}
We have $\det(A) = \prod_i \lambda_i$ (product over eigenvalues, a spectral invariant)
and $\det(A) = \prod_i s_i$ (product over SNF diagonal entries).
Hence $v_p(\det A) = \sum_i v_p(\lambda_i) = \sum_i v_p(s_i)$,
and the right-hand side is the same for any cospectral pair. \qed
\end{proof}

\begin{corollary}
For any cospectral pair $G_1, G_2$: $|\mathrm{Crit}(G_1)|_p = |\mathrm{Crit}(G_2)|_p$
for all primes $p$, where $|\cdot|_p$ denotes the $p$-part of the group order.
\end{corollary}

\subsection{The Spence Smith Ladder}

The Spence family consists of 28 non-isomorphic $\mathrm{SRG}(40,12,2,4)$, including
the symplectic polar space $W(3,3)$ and the orthogonal polar space $Q(4,3)$.
All share the adjacency spectrum $\{12^1, 2^{24}, (-4)^{15}\}$, hence by
Theorem~\ref{thm:transfer}, all have $v_2(\mathrm{Smith}) = v_2(\det A) = 56$.

The \emph{2-rank} of $G$ is $\rho_2(G) = \#\{i : 2 \mid s_i(G)\}$.
By GAP computation of all 28 SNFs \cite{Spence1975,GAP}, we obtain:

\begin{theorem}[Smith 2-Rank Ladder]
\label{thm:ladder}
The 28 Spence graphs partition by Smith 2-rank as:
$$\{17, 8, 2, 1\} \quad \text{(counts in each 2-rank class)}.$$
$W(3,3)$ and $Q(4,3)$ belong to the same 2-rank class and are
distinguished by their 2-Sylow subgroups of the critical group.
\end{theorem}

\begin{remark}
The Sunada-Gassmann structure of the pair $\{W(3,3), Q(4,3)\}$ (same Ihara
zeta, same Bartholdi zeta, same class number) is resolved at the level of
the Smith normal form: the critical groups differ in their 2-Sylow, providing
the finest group-theoretic non-isomorphism certificate.
\end{remark}

\subsection{Transcendence of the Smith Invariant}

\begin{theorem}[Smith Transcends Zeta]
\label{thm:transcend}
For the Spence family, the Smith critical group invariant is strictly finer
than any graph zeta function (Ihara, Hashimoto/edge, Bartholdi with any parameters).
\end{theorem}

\begin{proof}
By the Bass-Hashimoto theorem~\cite{Bass1992,Hashimoto1989}, for a $k$-regular
graph the Hashimoto (edge) zeta equals the Ihara zeta, both determined purely
by the adjacency spectrum. Since all 28 Spence graphs share the same spectrum,
they share the same zeta function. But the Smith 2-rank separates them into
4 classes, so the Smith invariant is strictly finer than any zeta invariant. \qed
\end{proof}

\begin{remark}
Theorem~\ref{thm:transcend} says that for this family, the first
\emph{non-spectral} invariant one encounters -- the Smith normal form -- is
already sufficient to separate 4 classes, while no spectral method separates even 2.
\end{remark}
